\section{Κεφάλαιο 1: Εισαγωγή (Introduction)}

\subsection{1.1 Το Κλινικό Πρόβλημα: Η Πρόκληση της Πρώιμης Διάγνωσης στην Ογκολογία}

Ο καρκίνος αποτελεί μία από τις κύριες αιτίες νοσηρότητας και θνησιμότητας παγκοσμίως, με τον Παγκόσμιο Οργανισμό Υγείας (ΠΟΥ) να καταγράφει σχεδόν 10 εκατομμύρια θανάτους ετησίως. Μεταξύ των διαφόρων τύπων κακοήθειας, ο Καρκίνος του Παχέος Εντέρου (Colorectal Cancer - CRC) κατέχει εξέχουσα θέση, αποτελώντας την τρίτη συχνότερη μορφή καρκίνου και τη δεύτερη αιτία θανάτου από καρκίνο παγκοσμίως [70]. Το παράδοξο του CRC είναι ότι, ενώ είναι εξαιρετικά θανατηφόρος στα προχωρημένα στάδια (μετάσταση), είναι ταυτόχρονα και μία από τις πιο ιάσιμες μορφές αν ανιχνευθεί εγκαίρως. Η φυσική ιστορία της νόσου ακολουθεί μια μακροχρόνια, σταδιακή εξέλιξη γνωστή ως \textbf{"Ακολουθία Αδενώματος-Καρκινώματος" (Adenoma-Carcinoma Sequence)}. Η μετάβαση από το φυσιολογικό επιθήλιο σε καλοήθη πολύποδα (αδένωμα) και τελικά σε διηθητικό καρκίνωμα μπορεί να διαρκέσει 10 έως 15 έτη [71]. Αυτό το χρονικό "παράθυρο ευκαιρίας" καθιστά τον CRC ιδανικό στόχο για προγράμματα προσυμπτωματικού ελέγχου (screening).

\subsubsection{1.1.1 Περιορισμοί των Συμβατικών Διαγνωστικών Μεθόδων}

Σήμερα, η διάγνωση βασίζεται κυρίως σε δύο πυλώνες:

\begin{enumerate}
\item \textbf{Επεμβατικές Μέθοδοι (Κολονοσκόπηση):} Αποτελεί το "Gold Standard" για τη διάγνωση. Ωστόσο, είναι μια επεμβατική διαδικασία, με υψηλό κόστος, δυσφορία για τον ασθενή και μικρό (αλλά υπαρκτό) κίνδυνο επιπλοκών (π.χ. διάτρηση εντέρου). Λόγω αυτών των παραγόντων, η συμμόρφωση του γενικού πληθυσμού στις οδηγίες screening παραμένει χαμηλή.
\item \textbf{Μη Επεμβατικές Μέθοδοι (FIT, CEA):} Τεστ όπως η ανίχνευση αιμοσφαιρίνης στα κοπράνα (FIT) έχουν χαμηλή ευαισθησία (sensitivity) για τα αρχικά στάδια (αδενώματα), ενώ οι κλασικοί καρκινικοί δείκτες στο αίμα (π.χ. CEA, CA19-9) αυξάνονται συνήθως όταν ο όγκος έχει ήδη μεγαλώσει σημαντικά, καθιστώντας τους ακατάλληλους για πρώιμη διάγνωση [74].
\end{enumerate}

Συνεπώς, υπάρχει μια επιτακτική κλινική ανάγκη για την ανάπτυξη νέων, ελάχιστα επεμβατικών και υψηλής ακρίβειας εργαλείων που μπορούν να ανιχνεύσουν τη νόσο στο στάδιο του αδενώματος.

\subsection{1.2 Η Επανάσταση της Υγρή Βιοψίας και της Πρωτεομικής}

Η λύση στο πρόβλημα της χαμηλής συμμόρφωσης έρχεται μέσω της \textbf{Υγρής Βιοψίας (Liquid Biopsy)}. Η έννοια βασίζεται στην παρατήρηση ότι οι όγκοι, ακόμα και σε μικροσκοπικό επίπεδο, αλληλεπιδρούν με το μικροπεριβάλλον τους και απελευθερώνουν βιομόρια (κύτταρα, DNA, εξωσώματα, πρωτεΐνες) στην κυκλοφορία του αίματος.

\subsubsection{1.2.1 Γιατί Πρωτεΐνες;}

Ενώ η γονιδιωματική (Genomics) αποκαλύπτει τι \emph{μπορεί} να συμβεί (βάσει των μεταλλάξεων), η Πρωτεομική (Proteomics) φανερώνει τι \emph{πραγματικά συμβαίνει} σε πραγματικό χρόνο. Οι πρωτεΐνες είναι οι τελικοί εκτελεστές των βιολογικών διεργασιών. Αλλαγές στη συγκέντρωση πρωτεϊνών που σχετίζονται με τη φλεγμονή, την αγγειογένεση (π.χ. VEGF) ή την αναδόμηση του ιστού (π.χ. MMPs) μπορούν να προηγηθούν της κλινικής εμφάνισης του όγκου κατά μήνες ή και χρόνια [75]. Η σύγχρονη τεχνολογία Φασματομετρίας Μάζας (Mass Spectrometry - MS) επιτρέπει τον ταυτόχρονο προσδιορισμό χιλιάδων πρωτεϊνών σε μια σταγόνα πλάσματος, δημιουργώντας ένα πολυδιάστατο "μοριακό αποτύπωμα" του ασθενούς.

\subsection{1.3 Το Πρόβλημα της Στατικής Ανάλυσης και η "Δυναμική" Λύση}

\subsubsection{1.3.1 Η Ανεπάρκεια του "Στιγμιότυπου" (Snapshot)}

Η πλειονότητα των διαγνωστικών μοντέλων Μηχανικής Μάθησης εκπαιδεύεται με στατικά δεδομένα: παίρνουμε μια μέτρηση $T_0$ και προσπαθούμε να ταξινομήσουμε τον ασθενή. Αυτή η προσέγγιση αγνοεί τη θεμελιώδη βιολογική αλήθεια πως η υγεία και η ασθένεια είναι δυναμικές διαδικασίες, όχι στατικές καταστάσεις.

Μια μεμονωμένη τιμή πρωτεΐνης μπορεί να είναι παραπλανητική λόγω της \textbf{Δια-ατομικής Μεταβλητότητας (Inter-patient variability)} αφού το "φυσιολογικό" επίπεδο μιας πρωτεΐνης διαφέρει από άνθρωπο σε άνθρωπο. Μια τιμή $X$ μπορεί να είναι φυσιολογική για τον ασθενή Α, αλλά παθολογική (υψηλή) για τον ασθενή Β, και από \textbf{Θορύβο (Noise)} με παράγοντες όπως ιώσεις, στρες ή διατροφή να επηρεάζουν προσωρινά τα επίπεδα πρωτεϊνών.

\subsubsection{1.3.2 Η Υπόθεση της Χρονικής Εξέλιξης (Research Hypothesis)}

Η παρούσα εργασία ασχολείται με την έννοια της \textbf{Διαχρονικής Πρωτεομικής (Longitudinal Proteomics)}. Η κεντρική υπόθεση εργασίας είναι:

\begin{quote}
\emph{"Η διαγνωστική πληροφορία μεγιστοποιείται όταν παρατηρούμε την τροχιά (trajectory) των βιοδεικτών στον χρόνο, αντί για την απόλυτη τιμή τους σε μια δεδομένη στιγμή."}
\end{quote}

Συγκεκριμένα, εξετάζουμε το σενάριο όπου λαμβάνονται δύο δείγματα ($T_0$ και $T_1$). Η υπόθεση είναι ότι η μετάβαση $T_0 - T_1$ αποκαλύπτει την "κλίση" της νόσου. Ένα σταθερό επίπεδο (ακόμα και ελαφρώς αυξημένο) μπορεί να υποδηλώνει καλοήθεια, ενώ μια αυξητική τάση μπορεί να σημάνει την έναρξη καρκινογένεσης. Για την αποκωδικοποίηση αυτών των χρονικών μοτίβων, οι παραδοσιακές στατιστικές μέθοδοι δεν επαρκούν και απαιτούνται αλγόριθμοι Βαθιάς Μάθησης (Deep Learning) ικανοί να μοντελοποιούν ακολουθίες.

\subsection{1.4 Σκοπός της Εργασίας}

Σκοπός της εργασίας είναι η ανάπτυξη ενός ολοκληρωμένου υπολογιστικού πλαισίου (computational framework) για την πρόβλεψη της κατάστασης του ασθενούς (Healthy / Adenoma / Cancer) βάσει χρονικών πρωτεομικών δεδομένων, καθώς και η ενασχόληση με το ερώτημα ποια οικογένεια αλγορίθμων είναι καταλληλότερη για τέτοιου είδους δεδομένα μικρού μεγέθους ($N < 100$). Η εργασία εκκινεί από την υπόθεση ότι οι αρχιτεκτονικές Βαθιάς Μάθησης (RNN/LSTM) είναι κατάλληλες για χρονοσειρές, και την εξετάζει συστηματικά μέσω συγκριτικών πειραμάτων.

Οι ειδικότεροι στόχοι περιλαμβάνουν:

\begin{enumerate}
\item \textbf{Συστηματική Σύγκριση Οικογενειών Αλγορίθμων:} Εμπειρική αξιολόγηση της απόδοσης Βαθιάς Μάθησης (Vanilla RNN, LSTM με Attention) έναντι Μεθόδων Συνόλου (Random Forest, Random Subspace, Rotation Forest, Custom Cascade Forest) και γραμμικών baseline μοντέλων (Logistic Regression) σε πραγματικό πρωτεομικό dataset.
\item \textbf{Ανάπτυξη Υβριδικών Αρχιτεκτονικών DL:} Σχεδιασμός και υλοποίηση μοντέλων που συνδυάζουν \textbf{1D-CNNs} (για εξαγωγή χαρακτηριστικών από το πρωτεϊνικό φάσμα), \textbf{Recurrent Neural Networks (RNNs/LSTMs)} (για τη μάθηση της χρονικής εξέλιξης), και \textbf{Attention Mechanisms} (για ερμηνευσιμότητα και ανακάλυψη βιοδεικτών).
\item \textbf{Σύγκριση Πειραματικών Παραδειγμάτων (Ablation Study):} Αυστηρή σύγκριση της προτεινόμενης χρονικής μεθόδου (Temporal) έναντι της καθιερωμένης στατικής (Static) και της απλοϊκής διαφορικής μεθόδου (Delta), για να ποσοτικοποιηθεί η αξία της χρονικής πληροφορίας.
\item \textbf{Αντιμετώπιση του "Small Data Problem":} Εφαρμογή εξειδικευμένων τεχνικών (Imputation, Regularization, Stratified Cross-Validation) για την εκπαίδευση μοντέλων σε datasets με περιορισμένο αριθμό ασθενών, ένα συχνό πρόβλημα στην ιατρική έρευνα.
\item \textbf{Αυστηρή Στατιστική Αξιολόγηση:} Πέρα από την απλή cross-validation, εφαρμογή Nested CV (αμερόληπτη εκτίμηση), Repeated CV, Bootstrap Confidence Intervals, και Permutation Tests, ώστε τα συμπεράσματα να διαθέτουν επιστημονική εγκυρότητα ανθεκτική στην κριτική για το μικρό μέγεθος δείγματος.
\item \textbf{Βιολογική Ερμηνευσιμότητα:} Ταυτοποίηση των πρωτεϊνών που οδηγούν τις αποφάσεις του μοντέλου μέσω univariate (ANOVA $F$-score), multivariate (Random Forest MDI) και attention-based αναλύσεων, προσφέροντας στους βιολόγους νέους υποψήφιους βιοδείκτες προς περαιτέρω διερεύνηση.
\end{enumerate}

Η εργασία ξεκίνησε με την προσδοκία ότι το Deep Learning θα κυριαρχούσε στα πειράματα, λόγω της εγγενούς ικανότητάς του να μοντελοποιεί χρονικές εξαρτήσεις. Η πειραματική πραγματικότητα, ωστόσο, οδήγησε στο συμπέρασμα της ότι σε small-N βιοϊατρικά datasets, τα απλούστερα μοντέλα (Logistic Regression) μπορούν να υπερτερούν.

\subsection{1.5 Δομή της Εργασίας}

Η εργασία διαρθρώνεται σε έξι κεφάλαια:

\begin{itemize}
\item \textbf{Κεφάλαιο 2:} Παρουσιάζει το θεωρητικό υπόβαθρο, καλύπτοντας τόσο τα Νευρωνικά Δίκτυα και τα Recurrent αρχιτεκτονικά (RNN, LSTM, Attention) όσο και τις Μεθόδους Συνόλου (Bagging, Random Forest, Rotation Forest, Deep Forest). Ανασκοπεί τη βιβλιογραφία σχετικά με την καταλληλότητα κάθε οικογένειας για datasets μικρού μεγέθους.
\item \textbf{Κεφάλαιο 3:} Περιγράφει αναλυτικά τη μεθοδολογία: προ-επεξεργασία πρωτεομικών δεδομένων, στρατηγικές imputation, feature selection, αρχιτεκτονικές των DL μοντέλων (Vanilla RNN, LSTM), υλοποίηση των ensemble μεθόδων, και τα τρία πειραματικά παραδείγματα (Static, Delta, Temporal). Καλύπτει επίσης τη βελτιστοποίηση υπερπαραμέτρων και την αυστηρή στατιστική αξιολόγηση (Nested CV, Permutation Test, Bootstrap CI).
\item \textbf{Κεφάλαιο 4:} Παραθέτει τα πειραματικά αποτελέσματα, συγκρίνοντας τις επιδόσεις των DL και ensemble μοντέλων, παρουσιάζοντας πίνακες σύγχυσης, στατιστικούς ελέγχους σημαντικότητας, learning curves, calibration analysis, και κλινικά σημαντικές μετρικές ανά κλάση.
\item \textbf{Κεφάλαιο 5:} Συζητά τα ευρήματα υπό το πρίσμα της βιολογίας και της μεθοδολογίας, αναλύοντας τον ρόλο των κορυφαίων πρωτεϊνών (MMP12, AREG, L1CAM) στην παθογένεια του καρκίνου, και ερμηνεύοντας γιατί τα γραμμικά μοντέλα υπερτέρησαν σε αυτή την περίπτωση.
\item \textbf{Κεφάλαιο 6:} Συνοψίζει τα συμπεράσματα, τα μεθοδολογικά διδάγματα, τους περιορισμούς, και προτείνει κατευθύνσεις για μελλοντική έρευνα.
\end{itemize}

